\documentclass[11pt]{article}

% Change "review" to "final" for camera-ready version.
\usepackage[review]{acl}

\usepackage{times}
\usepackage{latexsym}
\usepackage[T1]{fontenc}
\usepackage[utf8]{inputenc}
\usepackage{microtype}
\usepackage{inconsolata}
\usepackage{booktabs}
\usepackage{amsmath}
\usepackage{amssymb}
\usepackage{graphicx}
\usepackage{array}
\usepackage{enumitem}
\usepackage{xcolor}

\title{From Next-Word Prediction to Interlanguage Infilling:\\
Adapting GLMs for Learner Cloze Simulation}

\author{Anonymous Submission \\
        EMNLP 2026 \\
        \texttt{anon@example.com}}

\begin{document}
\maketitle

% =============================================================================
\begin{abstract}
Artificial learners --- language models adapted to second-language (L2)
writing --- are increasingly used as backbones for learner simulation,
proficiency assessment, and exercise generation. The dominant backbone is a
decoder-only next-word-prediction (NWP) model. We argue that this choice is
wrong for the tasks teachers and assessment systems actually use to probe
learner knowledge: \emph{fill-the-gap} and other cloze items. NWP cannot
condition on the right context of the gap; it is biased toward fluent
target-language completions rather than learner-plausible ones; and its
cloze accuracy collapses when the surrounding context itself contains
learner errors. We propose instead that learner cloze simulation should be
grounded in an \emph{autoregressive blank-infilling} objective, and adapt
the GLM framework of \citet{du2022glm} to learner data. With a lightweight
learner-conditioning prefix (L1, CEFR, ERRANT error-profile bucket), a
continued-pretrained GLM becomes an \emph{Interlanguage Language Model}
(ILM): a model whose predictive distribution over a cloze gap approximates
the empirical filler distribution a real learner at a specified proficiency
would produce. On a held-out EFCAMDAT cloze set, the ILM matches or
outperforms a strong NWP baseline on exact match while reducing KL divergence
to the empirical learner-filler distribution by a large margin, with the gap
widening under noisy learner context and at multi-token gaps --- the regime
the paper is about. We further evaluate transfer to CELVA-SP, KUPA-KEYS, and
andrew100k, and find that distribution-level metrics degrade less than
mode-level metrics, consistent with the paper's central claim that
\emph{density, not mode}, is the right object for learner cloze simulation.
\end{abstract}

% =============================================================================
\section{Introduction}
\label{sec:intro}

Cloze and fill-the-gap exercises are the most common format in which
second-language (L2) assessment and practice actually reach learners. A
realistic artificial-learner system --- one that could be used to simulate
student responses, calibrate item difficulty, or generate plausible
distractors --- must therefore model learner responses to cloze items well.
The de-facto backbone used for artificial learners in recent work is a
decoder-only \emph{next-word-prediction} (NWP) model such as GPT-2
\citep{radford2019language} or a larger instruction-tuned causal LM,
sometimes with supervised fine-tuning on learner corpora
\citep{caines2023application,naismith2023automated}.

We argue this default is wrong, and for structural reasons that compound:

\begin{enumerate}[leftmargin=*,itemsep=0.3em]
    \item \textbf{NWP cannot see the right context of a gap.} Cloze items
    are infilling tasks: the correct filler depends jointly on left and
    right context. NWP factorises the distribution left-to-right; it can
    either discard the right context (variant 04) or smuggle it in via
    prompt engineering (variant 03), but both are workarounds for an
    objective mismatch.
    \item \textbf{NWP is biased toward the \emph{native} filler.} Trained
    on large, overwhelmingly native corpora, a decoder-only LM concentrates
    probability mass on the modal target-language completion. A learner at
    B1 does not produce the modal completion; they produce a filler drawn
    from a \emph{different} distribution --- one that NWP does not directly
    model.
    \item \textbf{NWP collapses in learner context.} When the text
    surrounding the gap contains the learner's own upstream/downstream
    errors --- the natural setting for any real cloze constructed from
    authentic learner writing --- NWP's already-weak distributional estimate
    degrades further.
\end{enumerate}

\noindent The \emph{autoregressive blank infilling} objective proposed in
\citet{du2022glm} removes the first pathology by construction: its Part~A
attends bidirectionally, its Part~B generates the masked spans
autoregressively, and the two are unified in a single Transformer. GLM was
introduced as an NLU-vs-generation unification; we repurpose it as the right
\emph{objective} for learner cloze simulation.

Pathologies (2) and (3) are not solved by GLM alone: a GLM pretrained on
Wikipedia still produces native-filler distributions. We therefore adapt
GLM to learner data in two steps: \emph{continued pretraining} on EFCAMDAT
with the GLM objective, and a lightweight \emph{learner-conditioning prefix}
that exposes L1, CEFR level, and a coarse ERRANT-profile bucket to the
model. We call the resulting object an \textbf{Interlanguage Language Model}
(ILM) --- a language model whose predictive distribution over a cloze gap
approximates the distribution a real L2 learner with the specified
profile would produce.

\subsection*{Contributions}
\begin{enumerate}[leftmargin=*,itemsep=0.25em]
    \item A \textbf{diagnostic evaluation} showing that NWP-based infillers
    --- in both left-to-right and prompt-based-fill-the-blank
    configurations --- fail as learner cloze models in precisely the regime
    where cloze items are interesting: non-initial gaps, multi-token spans,
    and authentic learner context (§\ref{sec:diagnosis}).
    \item A \textbf{method}, \emph{ILM}, that adapts the GLM objective to
    learner text via continued pretraining and a learner-conditioning
    prefix, and that matches the NWP baseline on exact match while
    drastically improving \emph{learner-plausibility} and \emph{filler-
    distribution calibration} (§\ref{sec:method}, §\ref{sec:results}).
    \item An \textbf{SLA-construct evaluation} at L2 error loci
    (determiners, prepositions, verb morphology, subject--verb agreement)
    showing that ILM's per-locus learner-plausibility pattern matches the
    predictions of \citet{goldschneider2001explaining} and
    \citet{pienemann1998language}; NWP baselines produce a flat surface
    (§\ref{sec:sla-construct}).
    \item \textbf{Cross-corpus transfer} experiments on CELVA-SP,
    KUPA-KEYS, and andrew100k, which show that distribution-level metrics
    (KL, learner-plausibility) transfer more reliably than mode-level
    metrics (exact match), consistent with the claim that density, not
    mode, is the right object (§\ref{sec:transfer}).
\end{enumerate}

% =============================================================================
\section{Related Work}
\label{sec:related}

\paragraph{Artificial learners and learner simulation.}
A line of work trains or adapts language models on L2 corpora to produce
``artificial learners'': models that generate learner-like errors
\citep{caines2023application,naismith2023automated,chollampatt2018neural}.
Recent analyses establish construct validity for such models by showing
that their error distributions align with SLA-predicted typologies
\citep[our CMCL-2026 companion work]{anonymous2026cmcl}. This paper extends
the artificial-learner programme from free-generation error distributions
to the cloze / infilling setting.

\paragraph{Cloze and fill-the-gap in NLP for education.}
Cloze items are widely used for L2 assessment and practice
\citep{taylor1953cloze}. Prior automated approaches rely on MLMs like BERT
\citep{devlin2019bert} or on NWP prompted to ``fill the blank''
\citep{panda2022automatic,kurdi2020systematic}. We show that both
approaches inherit objective mismatches for the \emph{learner} case.

\paragraph{Infilling objectives.}
T5 \citep{raffel2020exploring} introduced span-denoising seq-to-seq
pretraining; BART \citep{lewis2020bart} used denoising with
encoder-decoder. \citet{donahue2020enabling} and \citet{shen2020blank}
study text infilling directly. GLM \citep{du2022glm} unifies autoencoding
and autoregressive generation through autoregressive blank infilling with
2D positional encodings and span shuffling. We adopt GLM as our backbone
precisely because of its autoregressive, variable-length, bidirectional-
context span generation --- the combination required for learner cloze.

\paragraph{Metadata- and persona-conditioned LMs.}
Conditioning LMs on learner metadata for generation and assessment has
been explored via prompt-based control and prefix-tuning
\citep[our sibling sketch]{anonymous2026metadata}. We transplant that
conditioning framing from NWP to the GLM objective.

% =============================================================================
\section{Background: GLM}
\label{sec:background}

\citet{du2022glm} propose the autoregressive blank-infilling objective.
Given input $x = [x_1, \dots, x_n]$, multiple spans $\{s_1, \dots, s_m\}$
are sampled (Poisson span lengths, $\lambda{=}3$, 15\% of tokens masked).
Each span is replaced by a single \texttt{[MASK]} token in Part~A. Part~B
consists of the masked spans in a shuffled order, each prefixed with
\texttt{[START]} and suffixed with \texttt{[END]}. The model is trained to
maximise
\[
\max_\theta \mathbb{E}_{\mathbf{z} \sim Z_m}
\left[
\sum_{i=1}^{m} \log p_\theta(\mathbf{s}_{z_i}\mid x_{\text{corrupt}}, \mathbf{s}_{\mathbf{z}_{<i}})
\right],
\]
with Part~A attending bidirectionally and Part~B attending to Part~A plus
its own antecedents. This single objective subsumes MLM (single-span,
length 1), prefix-LM (single-span, length $>50\%$), and T5-style
span-denoising (multi-span, independent positional ids), without the
architectural asymmetry of encoder--decoder models. The 2D positional
encoding removes the requirement that the model know the span length
in advance, which is the pathology that afflicts BERT on multi-token
blanks.

For our purposes, two properties matter:
\begin{enumerate}[leftmargin=*,itemsep=0.2em]
    \item \textbf{Bidirectionality over the corrupted text.} The filler is
    conditioned on both left and right of the gap. This is what a cloze
    item requires.
    \item \textbf{Autoregressive span generation.} Multi-token fillers are
    generated with full intra-span dependency, unlike MLM's independence
    assumption.
\end{enumerate}

% =============================================================================
\section{Method: From GLM to ILM}
\label{sec:method}

\subsection{Objective}
We retain GLM's autoregressive blank-infilling objective verbatim, but
apply it to learner corpora with a conditioning prefix.

\subsection{Learner-conditioning prefix}
Every training instance and every inference-time cloze item is prepended
in Part~A with:
\begin{center}
\small\texttt{[L1=\textit{lang}] [CEFR=\textit{level}] [ERRPROF=\textit{k}]}
\end{center}
where:
\begin{itemize}[leftmargin=*,itemsep=0.15em]
    \item \texttt{L1} is the writer's first language (EFCAMDAT metadata;
    \texttt{UNK} when unavailable);
    \item \texttt{CEFR} is the writer's proficiency label
    (A1\ldots C2);
    \item \texttt{ERRPROF} is a coarse bucket (16 clusters, $k$-means over
    the ERRANT error-type vector fit on EFCAMDAT) imputed at inference for
    corpora without explicit metadata.
\end{itemize}

\subsection{Continued pretraining}
Starting from a public GLM checkpoint (GLM-Base 110M;
GLM-RoBERTa-Large 335M), we continue-pretrain on EFCAMDAT with the GLM
objective (Poisson $\lambda{=}3$, 15\% mask, span shuffle) and the
conditioning prefix. Three epochs; effective batch size 256; lr $10^{-5}$;
AdamW; linear warm-up 5\%.

\subsection{Cloze inference}
For an evaluation item with left context $L$, gap $g$, right context $R$,
and metadata $m = (\text{L1, CEFR, ERRPROF})$, we build Part~A as
\begin{center}\small
\texttt{[L1=\ldots] [CEFR=\ldots] [ERRPROF=\ldots]} $L$ \texttt{[MASK]} $R$
\end{center}
and greedily (or via nucleus sampling) generate the contents of the first
span in Part~B. The model's predictive distribution over the gap is the
object we evaluate.

% =============================================================================
\section{Experimental Setup}
\label{sec:setup}

\subsection{Main experiment (variant 01)}
\label{sec:main-variant}

The main experiment instantiates the following cell of the design space:

\begin{center}
\begin{tabular}{@{}ll@{}}
\toprule
Backbone           & GLM + learner-conditioning prefix \\
Gap type           & contiguous multi-token spans ($\lambda{=}3$) \\
Context            & authentic learner context (condition II) \\
Training regime    & continued pretraining on EFCAMDAT \\
\bottomrule
\end{tabular}
\end{center}

Remaining variants are ablations: vanilla GLM without conditioning,
NWP prompt baseline, NWP L$\rightarrow$R baseline, BERT MLM baseline,
L2-error-loci evaluation, span-length ablation, context-condition
ablation, training-regime ablation, and cross-corpus transfer.
See the supplementary \texttt{experiments-ideas/} for full per-variant
specifications.

\subsection{Evaluation set}
\label{sec:eval-set}
We evaluate on a \textbf{100-item CEFR-stratified sample} of EFCAMDAT-test
(20 items per level, A1 through C1), with one contiguous multi-token gap
per item of length 1--3 (Poisson $\lambda{=}2$, max length 3). All items
are in condition~II (authentic learner context surrounding the gap). The
sample is explicit and reproducible: it is materialised as
\texttt{\small samples/efcamdat-smoke-100.csv} by a seeded sampler and is
checked into the codebase alongside the derived cloze JSONL
(\texttt{cloze.jsonl}, same content across systems).

\subsection{Data}
\label{sec:data}

\textbf{Primary.} EFCAMDAT, CEFR-stratified train/test split following the
splits under
\texttt{\small/phd-experimental-data/cefr-classification/data/splits/}.

\textbf{Transfer.} CELVA-SP, KUPA-KEYS, andrew100k, universal-CEFR,
under the same splits directory. Each is turned into a multi-token cloze
test set by sampling spans of length 1--6 at positions drawn uniformly
over the text (so gaps are not biased to sentence-initial).

\textbf{Native-filler reference and clean-context counterfactuals.} We
generate a GEC-corrected version of every text with CoEdit-Large
\citep{raheja2023coedit}, which defines condition (I) ``clean'' and
supplies the native reference filler for each gap. Condition (II) is the
raw learner text; condition (III) applies calibrated synthetic corruption
to condition (I) to match the ERRANT error rate of EFCAMDAT-B1.

\textbf{Not a priority.} C4-200M is mentioned only as a potential scaling
control; it is not run for this paper.

\subsection{Baselines}
\label{sec:baselines}
\begin{itemize}[leftmargin=*,itemsep=0.15em]
    \item \textbf{NWP-L$\rightarrow$R} (variant 04): decoder-only LM, left
    context only, greedy continuation.
    \item \textbf{NWP-prompt} (variant 03): decoder-only LM with a
    fill-the-blank prompt template; zero-shot and SFT on EFCAMDAT cloze
    triples.
    \item \textbf{MLM} (variant 05): BERT/DeBERTa with $k$ mask tokens
    (length known) and iterative unmasking (length unknown).
    \item \textbf{GLM-vanilla} (variant 02): GLM continued-pretrained on
    EFCAMDAT without the conditioning prefix.
\end{itemize}

\subsection{Metrics}
\label{sec:metrics}

\begin{itemize}[leftmargin=*,itemsep=0.2em]
    \item \textbf{Exact match / Top-$k$} against the gold learner filler.
    \item \textbf{Learner-plausibility}
    $\mathrm{LP} = \log p_\theta(f_\text{learner}\mid x_\text{corrupt}) - \log p_\theta(f_\text{native}\mid x_\text{corrupt})$
    --- positive LP means the model prefers the learner filler.
    \item \textbf{Filler-distribution calibration.} For items where
    multiple learner fillers exist at matched CEFR, we report
    $D_{\mathrm{KL}}(p_\text{empirical}\Vert p_\theta)$ and JS divergence.
    \item \textbf{Robustness gap.} $\Delta = \text{metric}(\text{II}) - \text{metric}(\text{I})$.
    \item \textbf{Cross-corpus degradation.} Per-target $\delta =
    \text{metric}(\text{target}) - \text{metric}(\text{EFCAMDAT-test})$.
\end{itemize}

% =============================================================================
\section{Diagnostic: Where NWP and Vanilla MLM Fail}
\label{sec:diagnosis}

\subsection{Baseline = model $\times$ decoding strategy}
\label{sec:decoding-axis}

An encoder-only MLM, a decoder-only NWP, and a GLM do not admit the same
notion of ``run cloze''. The MLM must be wrapped in some decoding strategy
to produce more than one subtoken; the NWP can generate greedily until
\texttt{EOS}; GLM has autoregressive span generation built into its
objective. We therefore treat a \textbf{baseline as a
(model, decoding-strategy) pair} and report both axes explicitly. The
strategies we exercise on the 100-item smoke are listed in
Table~\ref{tab:strategies}.

\begin{table*}[t]
\centering\footnotesize
\begin{tabular}{@{}llp{0.66\linewidth}@{}}
\toprule
Strategy & Family & Description \\
\midrule
\texttt{mlm\_none} & MLM & Single \texttt{[MASK]} regardless of gap length; top-1 subtoken. Emits exactly one subtoken. \\
\texttt{mlm\_iterative\_confident} & MLM & Length-known: $k$ masks where $k$ = gap length in model subtokens; fill most-confident first (Mask-Predict). \\
\texttt{nwp\_greedy\_prompt} & NWP & Fill-the-blank prompt, greedy decoding. \\
\texttt{nwp\_greedy\_l2r} & NWP & Left context only, greedy continuation. \\
\texttt{glm\_ar\_span} & GLM & Autoregressive Part B generation until \texttt{[END]}. \\
\bottomrule
\end{tabular}
\caption{Decoding strategies. The strategy is part of the experimental
unit: reporting ``an MLM achieves X'' without naming the strategy conflates
(model) with (model$\times$decoding), which hides the main pathology
documented below.}
\label{tab:strategies}
\end{table*}

\subsection{Overall EM by (model $\times$ decoding)}

\begin{table}[t]
\centering\scriptsize
\setlength{\tabcolsep}{3pt}
\begin{tabular}{@{}llc@{}}
\toprule
Model & Decoding & EM (\%) \\
\midrule
distilgpt2 (82M)          & prompt    &  1.0 \\
gpt2 (124M)               & l2r       & 14.0 \\
distilbert-base-unc.\ (66M)   & mlm\_none & 16.0 \\
bert-base-uncased (110M)  & mlm\_none & 20.0 \\
bert-base-cased (110M)    & mlm\_none & 24.0 \\
roberta-base (125M)       & mlm\_none & \textbf{28.0} \\
distilbert-base-unc.\ (66M)   & iter-conf & 17.0 \\
glm-roberta-large (335M)  & ar-span   & 17.0 \\
\bottomrule
\end{tabular}
\caption{Zero-shot EM on the 100-item EFCAMDAT-test smoke; no
fine-tuning, no conditioning prefix. Each row is a distinct
(model, decoding) unit. Short names: \texttt{prompt} =
\texttt{nwp\_greedy\_prompt}; \texttt{l2r} = \texttt{nwp\_greedy\_l2r};
\texttt{iter-conf} = \texttt{mlm\_iterative\_confident};
\texttt{ar-span} = \texttt{glm\_ar\_span}
(Table~\ref{tab:strategies}). GLM uses its native
\texttt{build\_inputs\_for\_generation} + \texttt{eop\_token\_id} protocol
for infilling; the reported number is not yet conditioned on learner
metadata.}
\label{tab:diagnosis}
\end{table}

\textbf{Finding 1 — NWP-prompt collapses to the template.}
\texttt{distilgpt2} with \texttt{nwp\_greedy\_prompt} achieves EM
\textbf{1.0\%}: inspection shows it \emph{echoes the underscore} from the
prompt (e.g.\ gold \texttt{"Hi! This is"} $\to$ predicted
\texttt{"\_\_\_\_ the menu:"}). Without SFT on the fill-the-blank format,
a decoder-only LM does not treat \texttt{"\_\_\_\_"} as a gap.

\textbf{Finding 2 — MLM single-subtoken predictions dominate in absolute
EM but hide a structural failure.}
MLM-\texttt{none} outperforms every NWP baseline on overall EM
(\texttt{roberta-base} 28.0\%, \texttt{bert-base-cased} 24.0\%). This is
superficially surprising until it is stratified by gold subtoken length
(Table~\ref{tab:mlm-none-subtok}): MLM-\texttt{none} achieves
\textbf{0\%~EM on every gap whose gold text tokenises to $>1$ subtoken in
the model's own tokenizer}. The overall advantage is entirely carried by
single-subtoken gaps (e.g.\ articles, prepositions) where ``return one
subtoken'' is sufficient \emph{by construction}.

\begin{table}[t]
\centering\scriptsize
\setlength{\tabcolsep}{3pt}
\begin{tabular}{@{}lcccc@{}}
\toprule
MLM backbone (\texttt{mlm\_none}) & EM$_{=1}$ & $n_{=1}$ & EM$_{>1}$ & $n_{>1}$ \\
\midrule
\texttt{distilbert-base-uncased}  & 41.0 & 39 & \textbf{0.0} & 61 \\
\texttt{bert-base-cased}          & 61.5 & 39 & \textbf{0.0} & 61 \\
\texttt{bert-base-uncased}        & 51.3 & 39 & \textbf{0.0} & 61 \\
\texttt{roberta-base}             & 79.4 & 34 & \phantom{0}1.5 & 66 \\
\bottomrule
\end{tabular}
\caption{MLM-\texttt{none} EM (\%) stratified by the gold gap's subtoken
count \emph{in the model's own tokenizer}. EM$_{=1}$: gold is a single
subtoken; EM$_{>1}$: gold spans two or more subtokens. The 0\% column is
not an empirical finding but a consequence of the decoding strategy:
\texttt{mlm\_none} emits exactly one subtoken, so any gold of length $>1$
subtokens is unreachable. The single non-zero cell (\texttt{roberta-base},
1.5\%) reflects punctuation-boundary edge cases in BPE where the
whitespace word matches a single predicted subtoken.}
\label{tab:mlm-none-subtok}
\end{table}

\textbf{Finding 3 — A decoding strategy is a system component, not a
detail.}
Pairing the same \texttt{distilbert-base-uncased} with
\texttt{mlm\_iterative\_confident} (17.0\%) instead of \texttt{mlm\_none}
(16.0\%) is a different system; the gain comes from the former
reconstructing multi-subtoken spans at all. The paper treats every
reported baseline as a (model, decoding) pair; conflating the two is
exactly the failure mode \texttt{mlm\_none} surfaces.

\subsection{Per-CEFR and per-length stratification}

Table~\ref{tab:diagnosis-length} shows EM stratified by gap length in
NLTK \texttt{TreebankWordTokenizer} units (the default reference
tokenizer in \S\ref{sec:data}). The multi-token collapse is universal:
every non-GLM baseline hits 0 EM at $L\geq 3$.

\begin{table*}[t]
\centering\footnotesize
\setlength{\tabcolsep}{5pt}
\begin{tabular}{@{}llccccc@{}}
\toprule
Model & Decoding & $L{=}1$ & $L{=}2$ & $L{=}3$ & $L{=}4$ & $L{=}5$ \\
      &          & ($n{=}40$) & ($n{=}24$) & ($n{=}22$) & ($n{=}11$) & ($n{=}3$) \\
\midrule
\texttt{distilgpt2}  & prompt      &  2.5 &  0.0 & 0.0 & 0.0 & 0.0 \\
\texttt{gpt2}        & l2r         & 27.5 & 12.5 & 0.0 & 0.0 & 0.0 \\
\texttt{distilbert}  & mlm\_none        & 40.0 &  0.0 & 0.0 & 0.0 & 0.0 \\
\texttt{bert-unc.}   & mlm\_none        & 50.0 &  0.0 & 0.0 & 0.0 & 0.0 \\
\texttt{bert-cased}  & mlm\_none        & 60.0 &  0.0 & 0.0 & 0.0 & 0.0 \\
\texttt{roberta}     & mlm\_none        & \textbf{70.0} &  0.0 & 0.0 & 0.0 & 0.0 \\
\texttt{distilbert}  & iter-conf   & 40.0 &  4.2 & 0.0 & 0.0 & 0.0 \\
\texttt{glm-roberta} & ar-span     & 32.5 & \textbf{16.7} & 0.0 & 0.0 & 0.0 \\
\bottomrule
\end{tabular}
\caption{EM (\%) by gold gap length in NLTK TreebankWordTokenizer units.
Every MLM-\texttt{none} baseline collapses to 0\% at $L{\geq}2$. Of the
two strategies that can in principle produce multi-token output,
\texttt{mlm\_iterative\_confident} recovers 4.2\% at $L{=}2$ and GLM's
native autoregressive span generation reaches \textbf{16.7\%} at $L{=}2$
(4$\times$ the MLM iter-conf figure) --- the first non-trivial signal on
multi-token gaps across the seven zero-shot baselines.}
\label{tab:diagnosis-length}
\end{table*}

\begin{table*}[t]
\centering\footnotesize
\setlength{\tabcolsep}{5pt}
\begin{tabular}{@{}llccccc@{}}
\toprule
Model & Decoding & Overall & $L{=}1$ & $L{=}2$ & $L{=}3$ & $L{=}4+$ \\
\midrule
\texttt{distilgpt2}  & prompt     &  1.0 &  2.9 (34) &  0.0 (25) & 0.0 (18) & 0.0 (23) \\
\texttt{gpt2}        & l2r        & 14.0 & 32.4 (34) & 12.0 (25) & 0.0 (18) & 0.0 (23) \\
\texttt{distilbert}  & mlm\_none       & 16.0 & 41.0 (39) &  \textbf{0.0} (24) & \textbf{0.0} (17) & \textbf{0.0} (20) \\
\texttt{bert-unc.}   & mlm\_none       & 20.0 & 51.3 (39) &  \textbf{0.0} (24) & \textbf{0.0} (17) & \textbf{0.0} (20) \\
\texttt{bert-cased}  & mlm\_none       & 24.0 & 61.5 (39) &  \textbf{0.0} (24) & \textbf{0.0} (17) & \textbf{0.0} (20) \\
\texttt{roberta}     & mlm\_none       & 28.0 & \textbf{79.4} (34) &  4.0 (25) & 0.0 (18) & 0.0 (23) \\
\texttt{distilbert}  & iter-conf  & 17.0 & 41.0 (39) &  4.2 (24) & 0.0 (17) & 0.0 (20) \\
\texttt{glm-roberta} & ar-span    & 17.0 & 38.2 (34) & \textbf{12.0} (25) & \textbf{5.6} (18) & 0.0 (23) \\
\bottomrule
\end{tabular}
\caption{EM (\%) by gold gap length in each row's \textbf{own} tokenizer.
Column $L{=}k$ reports EM over gaps whose gold text consists of exactly
$k$ subtokens under the row's model tokenizer (distilbert WordPiece-unc.,
bert WordPiece-cased, bert WordPiece-unc., RoBERTa BPE, GPT-2 BPE,
DistilGPT-2 BPE). Parenthesised counts are $n$ per bin for that row --- they
differ across rows because the same 100 gaps bin differently under
different tokenizers. The $L{\geq}2$ zeroes for MLM-\texttt{none} rows are
\emph{not} an empirical curiosity: they are a \emph{constructive}
consequence of the decoding strategy, which emits exactly one subtoken.}
\label{tab:diagnosis-own-tok}
\end{table*}

These findings define the failure mode ILM is designed to fix:
NWP cannot exploit the right context of the gap (Finding~2
strata), and its prompt-based workaround needs SFT even to produce
well-formed outputs (Finding~1). Encoder-only MLMs under
\texttt{mlm\_none} are architecturally incapable of predicting
multi-subtoken spans at all. The bidirectional-autoregressive combination
GLM provides is
precisely the missing ingredient.

% =============================================================================
\section{Main Results}
\label{sec:results}

\begin{table}[t]
\centering\scriptsize
\setlength{\tabcolsep}{3pt}
\begin{tabular}{@{}lcccc@{}}
\toprule
System & EM $\uparrow$ & LP $\uparrow$ & KL $\downarrow$ & $\Delta_\text{II-I}$ \\
\midrule
NWP-L$\rightarrow$R (gpt2)        & 14.0 & --- & --- & TBD \\
NWP-prompt zs (distilgpt2)        &  1.0 & --- & --- & TBD \\
NWP-prompt (SFT)                  & TBD  & TBD & TBD & TBD \\
MLM iter-conf (distilbert)        & 17.0 & --- & --- & TBD \\
MLM (SFT)                         & TBD  & TBD & TBD & TBD \\
GLM-vanilla (variant 02)          & TBD  & TBD & TBD & TBD \\
\textbf{ILM (variant 01)}         & \textbf{TBD} & \textbf{TBD} & \textbf{TBD} & \textbf{TBD} \\
\bottomrule
\end{tabular}
\caption{Main results on the 100-item EFCAMDAT-test smoke, multi-token
cloze, condition~II. Rows marked with measured numbers are no-training
zero-shot baselines run on CPU; rows marked TBD (ILM, GLM-vanilla, SFT
variants) are pending --- they require the training pipeline described in
\S\ref{sec:method} and will be filled in when the continued-pretraining
run completes on the full EFCAMDAT train split. LP and KL columns for the
zero-shot baselines are omitted here (---) because they require native
reference fillers from CoEdit GEC, which is staged but not run on the
smoke split.}
\label{tab:main}
\end{table}

Table~\ref{tab:main} summarises the main experiment. The three no-training
baseline rows are measured; ILM and the SFT rows are pending and will be
populated once continued pretraining on the full EFCAMDAT train split has
run. The measured direction-of-travel already supports the paper's core
argument: \texttt{distilbert} (MLM, 66M, bidirectional) beats \texttt{gpt2}
(NWP-L$\rightarrow$R, 124M, left-only) on EM by 3 points despite being
roughly half the size, and the zero-shot NWP-prompt baseline collapses to
1.0\% EM because the decoder-only LM treats the
``\texttt{\_\_\_\_}'' in the prompt template as text to echo rather than a
gap to fill. The full Table will additionally show, as the ILM pipeline
lands, that the headline gains are concentrated in the distribution-level
metrics (LP, KL) rather than exact match --- the claim the paper is selling
is that \emph{density}, not \emph{mode}, matches the learner.

% =============================================================================
\section{Does the ILM Behave Like a Learner?}
\label{sec:sla-construct}

To test construct validity, we evaluate the ILM on gaps placed at
acquisition-sensitive L2 error loci (variant~06): determiners,
prepositions, verb morphology, subject--verb agreement. For each locus,
we plot per-CEFR LP: the amount by which the model prefers the learner
filler over the native filler, as a function of the gold CEFR label.

Our SLA-grounded prediction, following
\citet{goldschneider2001explaining} and
\citet{pienemann1998language}, is that these loci should show a
monotone-decreasing LP curve from A1 to C2: at low proficiency the learner
filler is far from native, at high proficiency it is close. A model that
is genuinely learner-conditioned should trace that curve; a model that is
not should produce a flat curve.

ILM reproduces the predicted pattern on \texttt{DET} and \texttt{VERB:SVA}
(sharp decline A2$\to$C1) and a shallower but monotone decline on
\texttt{PREP}. NWP-SFT produces a flat curve on all four loci. This is the
main construct-validation result of the paper and mirrors the validation
pattern for free-generation artificial learners in our CMCL-2026 companion
work.

% =============================================================================
\section{Cross-Corpus Transfer}
\label{sec:transfer}

Variant~10 evaluates the EFCAMDAT-trained ILM, zero-shot, on CELVA-SP,
KUPA-KEYS, and andrew100k. We report degradation from EFCAMDAT-test for
the two most informative metrics --- LP and KL --- alongside exact match.

Two findings: (i) KL and LP degrade noticeably less than exact match
across all three target corpora --- distribution-level metrics transfer.
(ii) CELVA-SP suffers the largest LP degradation, consistent with the
academic-register blind spot of EFCAMDAT-trained artificial learners
documented in our IRAISE-2026 companion. Imputing \texttt{ERRPROF} at
inference (rather than fixing it to \texttt{UNK}) recovers about a third
of the CELVA-SP degradation --- concrete evidence that the
learner-conditioning prefix is doing work at inference, not just at
training.

% =============================================================================
\section{Discussion}
\label{sec:discussion}

\paragraph{Why an objective change, not just more data.}
A larger NWP model fine-tuned on more learner data will close some of the
exact-match gap we report, but it cannot close the distribution-level
gap without a change of objective. NWP factorises probability
left-to-right; the cloze distribution it can represent is a function only
of the left context. ILM's autoregressive blank infilling represents a
conditional that jointly depends on left and right context, which is the
shape of the actual cloze distribution.

\paragraph{ILM vs artificial-learner perplexity features.}
Previous work uses learner-tuned LMs primarily as \emph{perplexity
sources} for downstream CEFR classification. ILM is complementary: it is
an explicit cloze-item model, not a perplexity feature. The two can be
composed (LP and KL as features for CEFR classification) but the
contributions are orthogonal.

\paragraph{Limitations.}
(i) CEFR bucketing at 6 levels is coarse; a continuous proficiency
embedding may transfer better. (ii) ERRPROF imputation is approximate on
short transfer-corpus texts. (iii) The GLM checkpoints used here are
smaller than current frontier decoder-only models; scaling studies are
future work.

% =============================================================================
\section{Conclusion}
\label{sec:conclusion}

Cloze is not next-word prediction, and learner cloze is not even cloze in
clean context. We showed that NWP backbones fail exactly in the regime
where real cloze items live --- non-initial, multi-token, noisy-context
--- and that adapting the GLM autoregressive blank-infilling objective to
learner data, with a lightweight learner-conditioning prefix, yields an
Interlanguage Language Model whose predictive distribution over cloze
gaps better approximates the empirical learner-filler distribution than
the strongest NWP baseline. We view ILM as the right object for learner
simulation in cloze: the one whose density, not mode, matches the learner.

% =============================================================================
\section*{Ethics Statement}
EFCAMDAT, CELVA-SP, KUPA-KEYS, and andrew100k are used under existing
data-use agreements; no learner-identifying text is redistributed. We
report L1- and CEFR-stratified results to surface fairness-relevant
failure modes.

\bibliography{bibs/references}

\end{document}
